\subsubsection{Measurement Operators and Projectors}

In this section, we will explain \textbf{how measurement is actually performed mathematically}.

\highspace
\begin{flushleft}
    \textcolor{Green3}{\faIcon{book} \textbf{Measurement Operators}}
\end{flushleft}
After introducing the Spectral Theorem, we know that any observable can be written as:
\begin{equation*}
    O = \sum_k \lambda_k \proj{k}
\end{equation*}
Where:
\begin{itemize}
    \item $\lambda_k$ are the possible measurement outcomes (eigenvalues);
    \item $\ket{k}$ are the corresponding eigenvectors;
    \item $\proj{k}$ are projection operators onto the eigenvectors.
\end{itemize}
At this point, a natural question arises: \hl{what role do these projectors play in the measurement process?}

\highspace
The answer is that each projector:
\begin{equation*}
    M_k = \proj{k}
\end{equation*}
Is called a \definition{Measurement Operator} (\autopageref{eq:measurement-operator}). It represents the mathematical operation associated with the measurement outcome $\lambda_k$.

\begin{definitionbox}[: Measurement Operator]
    Given an observable (\autopageref{sec:observables-what-are-we-measuring}):
    \begin{equation*}
        O = \sum_k \lambda_k \proj{k}
    \end{equation*}
    The \definition{Measurement Operator} associated with the outcome $\lambda_k$ is:
    \begin{equation*}
        M_k = \proj{k} = \begin{pmatrix}
            k_{1} \\
            k_{2}
        \end{pmatrix}
        \begin{pmatrix}
            \overline{k_{1}} & \overline{k_{2}}
        \end{pmatrix}
        =
        \begin{pmatrix}
            k_{1} \overline{k_{1}} & k_{1} \overline{k_{2}} \\
            k_{2} \overline{k_{1}} & k_{2} \overline{k_{2}}
        \end{pmatrix}
    \end{equation*}
\end{definitionbox}

\noindent
In other words, \hl{each possible measurement outcome has an associated projector} that:
\begin{enumerate}
    \item \textbf{Extracts the component} of the quantum state along the eigenvector $\ket{k}$;
    \item \textbf{Determines the probability} of obtaining the corresponding eigenvalue $\lambda_k$;
    \item \textbf{Produces the post-measurement} state after collapse.
\end{enumerate}
For this reason, a measurement operator can be thought of as a \textbf{detector} tuned to a specific eigenstate. If the detector corresponding to $\ket{k}$ ``clicks'', then:
\begin{itemize}
    \item The measurement outcome is $\lambda_k$;
    \item The system collapses to the state $\ket{k}$.
\end{itemize}
For example, suppose an observable has two eigenstates \ket{0} and \ket{1}. Then there are two measurement operators:
\begin{equation*}
    M_0 = \proj{0},
    \qquad
    M_1 = \proj{1}
\end{equation*}
We can imagine them as two detectors:
\begin{itemize}
    \item Detector $M_0$ looks for the component of the state along \ket{0}.
    \item Detector $M_1$ looks for the component of the state along \ket{1}.
\end{itemize}
Suppose the state is:
\begin{equation*}
    \ket{\psi} = \alpha\ket{0} + \beta\ket{1}
\end{equation*}
This means that part of the state points along \ket{0}, and part of the state points along \ket{1}. The detector $M_0 = \proj{0}$ asks ``\emph{how much of the state is aligned with \ket{0}}?'' and extracts $\alpha\ket{0}$. The detector $M_1 = \proj{1}$ asks ``\emph{how much of the state is aligned with \ket{1}}?'' and extracts $\beta\ket{1}$.

\begin{examplebox}[: Example: Measurement Operators for Pauli-Z]
    The Pauli-Z observable can be written as:
    \begin{equation*}
        Z = \proj{0} - \proj{1}
    \end{equation*}

    Therefore, the associated measurement operators are:
    \begin{equation*}
        M_0 = \proj{0},
        \qquad
        M_1 = \proj{1}
    \end{equation*}

    These operators correspond to the two possible outcomes:
    \begin{itemize}
        \item $+1$, associated with the eigenstate $\ket{0}$;
        \item $-1$, associated with the eigenstate $\ket{1}$.
    \end{itemize}
\end{examplebox}

\noindent
Suppose the system is in the state:
\begin{equation*}
    \ket{\psi}
    =
    \alpha \ket{0}
    +
    \beta \ket{1}
\end{equation*}
Applying the measurement operator $M_0$:
\begin{equation*}
    M_0 \ket{\psi}
    =
    \ket{0}\bra{0}
    \left(
        \alpha \ket{0}
        +
        \beta \ket{1}
    \right)
    =
    \alpha \ket{0}
\end{equation*}
This means that the operator $M_0$ extracts only the component of the state aligned with $\ket{0}$ and removes the orthogonal component $\ket{1}$.

\highspace
The probability of obtaining the corresponding outcome is:
\begin{equation*}
    p_0
    =
    \bra{\psi} M_0 \ket{\psi}
    =
    |\alpha|^2
\end{equation*}
Similarly:
\begin{equation*}
    p_1
    =
    \bra{\psi} M_1 \ket{\psi}
    =
    |\beta|^2
\end{equation*}
Therefore, \textbf{measurement operators} provide the complete mathematical description of the measurement process: they \hl{isolate the relevant component} of the quantum state, \hl{determine the probability} of each outcome, and \hl{specify the state} after the measurement.

\highspace
\begin{flushleft}
    \textcolor{Green3}{\faIcon{book} \textbf{Completeness of the Basis}}
\end{flushleft}
In the previous example, we defined two measurement operators $M_0$ and $M_1$ corresponding to the two eigenstates of the Pauli-Z observable. But what if we had an observable with more eigenstates? Or what if we had a different observable with a completely different set of eigenvectors? Therefore, \hl{how do we know that the measurement operators we defined account for all possible outcomes of the measurement?}

\highspace
The answer is that the eigenvectors of a Hermitian observable form a \textbf{complete orthonormal basis}. This means that any quantum state can be expressed as a linear combination of these eigenvectors:
\begin{equation*}
    \ket{\psi} = \sum_k c_k \ket{k}
\end{equation*}
Mathematically, this property is expressed by the \definition{Completeness Relation}:
\begin{equation*}
    \sum_k \proj{k} = I
\end{equation*}
Where $I$ is the identity operator (\autopageref{sec:identity-gate}).

\begin{definitionbox}[: Completeness Relation]
    \hl{If $\{\ket{k}\}$ is a complete orthonormal basis}, then:
    \begin{equation}\label{eq:completeness-relation}
        \sum_k \proj{k} = I
    \end{equation}
    Since the measurement operators are defined as:
    \begin{equation*}
        M_k = \proj{k}
    \end{equation*}
    It follows that:
    \begin{equation*}
        \sum_k M_k = I
    \end{equation*}
    This means that the \hl{set of measurement operators is complete and accounts for all possible outcomes of the measurement.}
\end{definitionbox}

\begin{deepeningbox}[: When a set of vectors $\left\{\ket{k}\right\}$ is a complete orthonormal basis?]
    A set of vectors $\left\{\ket{k}\right\}$ is called a \definition{Complete Orthonormal Basis} when it satisfies \textbf{three conditions}:
    \begin{enumerate}
        \item \textbf{Normalized}, each vector has length 1: $\braket{k}{k} = 1$.
        \item \textbf{Orthogonal}, different vectors are perpendicular: $\braket{k}{j} = 0$ for $k \neq j$.
        \item \textbf{Complete}, the vectors span the entire Hilbert space, meaning that every quantum state can be written as a linear combination of them. That is, for any state \ket{\psi}, there exist coefficients $c_k$ such that:
        \begin{equation*}
            \ket{\psi} = \sum_k c_k \ket{k}
        \end{equation*}
    \end{enumerate}
    Compactly, the \emph{normalized} and \emph{orthogonal} conditions are written as:
    \begin{equation}\label{eq:orthonormality-relation}
        \braket{k}{j} = \delta_{kj}
    \end{equation}
    Where $\delta_{kj}$ is the \definition{Kronecker delta}, which is 1 if $k = j$ and 0 otherwise:
    \begin{equation}\label{eq:kronecker-delta}
        \delta_{kj} =
        \begin{cases}
            1 & k = j \\
            0 & k \neq j
        \end{cases}
    \end{equation}
    And the \emph{completeness} condition is expressed by the completeness relation:
    \begin{equation*}
        \sum_k \proj{k} = I
    \end{equation*}
    \begin{examplebox}[: Single-Qubit Computational Basis]
        For a single qubit, the computational basis consists of the vectors:
        \begin{equation*}
            \ket{0} =
            \begin{bmatrix}
                1 \\
                0
            \end{bmatrix},
            \qquad
            \ket{1} =
            \begin{bmatrix}
                0 \\
                1
            \end{bmatrix}
        \end{equation*}
        These vectors are:
        \begin{itemize}
            \item Normalized: $\braket{0}{0} = 1$ and $\braket{1}{1} = 1$.
            \item Orthogonal: $\braket{0}{1} = 0$.
            \item Complete: any single-qubit state can be expressed as a linear combination of $\ket{0}$ and $\ket{1}$.
            \begin{equation*}
                \ket{\psi} = \alpha \ket{0} + \beta \ket{1}
            \end{equation*}
        \end{itemize}
    \end{examplebox}
\end{deepeningbox}

\noindent
This \autoref{eq:completeness-relation} states that the \textbf{set of measurement operators is complete}: together, they \hl{describe every possible outcome of the measurement.}

\highspace
\textcolor{Green3}{\faIcon{question-circle} \textbf{Why is the completeness relation important?}}
To understand why this relation is important, let us apply both sides to an arbitrary state:
\begin{equation*}
    \ket{\psi}
    =
    \sum_k c_k \ket{k}
\end{equation*}
Where $c_k$ are the coefficients of $\ket{\psi}$ in the basis $\{\ket{k}\}$.

\highspace
Now take the inner product with $\bra{k}$ on both sides:
\begin{equation*}
    \braket{k}{\psi}
    =
    \bra{k}
    \left(
        \sum_j c_j \ket{j}
    \right)
\end{equation*}
Using linearity of the inner product:
\begin{equation*}
    \braket{k}{\psi}
    =
    \sum_j c_j \braket{k}{j}
\end{equation*}
Since the basis is orthonormal, $\braket{k}{j} = \delta_{kj}$, so only the term where $j = k$ survives (\autoref{eq:orthonormality-relation}, \autopageref{eq:orthonormality-relation}):
\begin{equation*}
    \braket{k}{j}
    =
    \delta_{kj}
    =
    \begin{cases}
        1 & k = j \\
        0 & k \neq j
    \end{cases}
\end{equation*}
Thanks to the Kronecker delta, all terms in the sum vanish except the one where $j = k$:
\begin{equation*}
    \braket{k}{\psi}
    =
    c_k \braket{k}{k}
    =
    c_k \cdot 1
    =
    c_k
\end{equation*}
In other words, the bra $\bra{k}$ acts as a coefficient extractor: it selects exactly the amount of $\ket{\psi}$ aligned with the basis vector $\ket{k}$.

\highspace
Finally, apply the sum of all projectors to $\ket{\psi}$:
\begin{equation*}
    \left(\sum_k \proj{k}\right)\ket{\psi}
    =
    \sum_k \ket{k}\braket{k}{\psi}
    =
    \sum_k \ket{k} c_k
    =
    \sum_k c_k \ket{k}
    =
    \ket{\psi}
\end{equation*}
This shows that the \textbf{sum of all projectors reproduces the original state exactly}. In other words, \hl{each projector $\proj{k}$ extracts one component of the state, and when all these components are added together, the original state is reconstructed without losing any information.} Therefore:
\begin{equation*}
    \sum_k \proj{k} = I
\end{equation*}

\begin{examplebox}[: Computational Basis]
    For a single qubit, the computational basis is:
    \begin{equation*}
        \ket{0},
        \qquad
        \ket{1}
    \end{equation*}

    The corresponding projectors are:
    \begin{equation*}
        \proj{0}
        =
        \begin{bmatrix}
            1 & 0 \\
            0 & 0
        \end{bmatrix},
        \qquad
        \proj{1}
        =
        \begin{bmatrix}
            0 & 0 \\
            0 & 1
        \end{bmatrix}
    \end{equation*}

    Their sum is:
    \begin{equation*}
        \proj{0}
        +
        \proj{1}
        =
        \begin{bmatrix}
            1 & 0 \\
            0 & 0
        \end{bmatrix}
        +
        \begin{bmatrix}
            0 & 0 \\
            0 & 1
        \end{bmatrix}
        =
        \begin{bmatrix}
            1 & 0 \\
            0 & 1
        \end{bmatrix}
        =
        I
    \end{equation*}
\end{examplebox}

\noindent
The completeness relation also guarantees that the probabilities of all possible outcomes sum to one. If:
\begin{equation*}
    p_k
    =
    \bra{\psi} M_k \ket{\psi}
\end{equation*}
Then:
\begin{align*}
    \sum_k p_k
    &=
    \sum_k \bra{\psi} M_k \ket{\psi}
    \\
    &=
    \bra{\psi}
    \left(
        \sum_k M_k
    \right)
    \ket{\psi}
    \\
    &=
    \bra{\psi} I \ket{\psi}
    \\
    &=
    \braket{\psi}{\psi}
    \\
    &= 1
\end{align*}
Therefore, the \textbf{completeness relation} ensures that:
\begin{itemize}
    \item \textbf{All} possible \textbf{measurement outcomes} are included;
    \item \textbf{No probability is lost};
    \item The \textbf{total probability} is always \textbf{equal to one}.
\end{itemize}